% =====================================================================
% MedXplain - IEEE conference paper, FIRST DRAFT
%
% Compile:
%   pdflatex medxplain_ieee.tex   (run twice for references)
%
% Needs IEEEtran.cls - ships with TeX Live/MiKTeX, or download from
% https://www.ieee.org/conferences/publishing/templates.html
%
% DRAFT STATUS: every number that has not been measured yet renders in red
% via \TBD or \ph{}. Set \draftnumsfalse below to turn the colouring off once
% the real results are in. Do not remove the macros - they are the only thing
% stopping a placeholder from being mistaken for a result.
% =====================================================================

\documentclass[conference]{IEEEtran}
\IEEEoverridecommandlockouts

\usepackage{cite}
\usepackage{amsmath,amssymb}
\usepackage{graphicx}
\usepackage{booktabs}
\usepackage{xcolor}
\usepackage{url}

\usepackage[hidelinks]{hyperref}

% --- draft placeholder machinery -------------------------------------
\newif\ifdraftnums\draftnumstrue   % <-- \draftnumsfalse when results are final
\newcommand{\TBD}{\ifdraftnums\textcolor{red}{\textbf{[TBD]}}\else--\fi}
\newcommand{\ph}[1]{\ifdraftnums\textcolor{red}{#1}\else#1\fi}
\newcommand{\note}[1]{\ifdraftnums\textcolor{blue}{\small[#1]}\fi}
% ---------------------------------------------------------------------

\begin{document}

\title{Faithful, Stable, or Plausible? A Three-Axis Evaluation of\\
Explanation Methods for Transformer-Based Clinical Risk Prediction}

\author{\IEEEauthorblockN{Anwesha Dash}
\IEEEauthorblockA{\textit{\ph{Department}} \\
\textit{\ph{Institution}}\\
\ph{City, Country} \\
anwesha.dash3010@gmail.com}
\and
\IEEEauthorblockN{\ph{Co-author / Supervisor}}
\IEEEauthorblockA{\textit{\ph{Department}} \\
\textit{\ph{Institution}}\\
\ph{City, Country} \\
\ph{email}}
}

\maketitle

\begin{abstract}
Transformer language models pretrained on clinical text predict patient outcomes
from unstructured notes with high accuracy, but clinicians will not act on an
unexplained risk score. Token-level attribution methods are routinely added to
close this gap, and are almost always validated by visual inspection of a single
method. We argue that clinical deployment requires three distinct properties of
an explanation --- \emph{faithfulness} to the model's actual decision process,
\emph{stability} under the paraphrase and abbreviation variance endemic to
clinical documentation, and \emph{plausibility} to a clinician, including the
direction of each contribution --- and that the field has no evidence these
properties are satisfied together. We evaluate six explanation methods spanning
attention-based, perturbation-based and gradient-based families on a single
fine-tuned clinical transformer over MIMIC-III discharge summaries, measuring all
three axes plus computational cost on the same model and the same notes. We find
that \TBD. We additionally propose a rank-ensemble explainer that recovers \TBD\
of SHAP's faithfulness at \TBD\ of its computational cost. Our results indicate
that the choice of explanation method is a documented trade-off with clinical
consequences rather than an implementation detail. Code and the full evaluation
harness are released.
\end{abstract}

\begin{IEEEkeywords}
clinical natural language processing, explainable AI, interpretability,
faithfulness, transformers, risk prediction, MIMIC-III
\end{IEEEkeywords}

% =====================================================================
\section{Introduction}
% =====================================================================

Clinical notes contain the richest information in the medical record. Severity,
trajectory, response to treatment and goals of care are documented in narrative
prose and are largely absent from the structured fields --- laboratory values,
vital signs, billing codes --- on which conventional risk models are built.
Transformer language models pretrained on clinical or biomedical corpora, such as
ClinicalBERT~\cite{alsentzer2019}, BioBERT~\cite{lee2020} and
PubMedBERT~\cite{gu2021}, extract this signal effectively, and clinical outcome
prediction from free text is by now a well-established
task~\cite{huang2019,vanaken2021}.

These models nevertheless output a bare probability. A clinician presented with
``\,0.71 risk of in-hospital mortality\,'' and no justification has no basis on
which to accept, override, or learn from the prediction, and a growing body of
regulatory and professional guidance treats explainability as a precondition for
clinical decision support rather than a desirable extra.

The standard response is to attach a token-level attribution method --- most
often attention weights, sometimes SHAP~\cite{lundberg2017} --- and to render the
note with the highest-scoring words highlighted. The resulting figure is then
validated by inspection: the highlighted phrases look clinically sensible, and
the system is declared interpretable. This is a weak standard, for a reason that
is well understood in the broader NLP literature: an explanation can be entirely
plausible to a human while being a poor description of what the model actually
did~\cite{jain2019,wiegreffe2019}. The ERASER benchmark~\cite{deyoung2020}
introduced quantitative \emph{faithfulness} metrics precisely to address this,
but they have seen limited uptake in clinical NLP.

We argue the problem is not only that faithfulness goes unmeasured, but that
clinical deployment imposes \emph{three} distinct requirements that the
literature evaluates separately, if at all:

\begin{enumerate}
\item \textbf{Faithfulness.} Do the highlighted tokens actually drive the
      prediction, or merely co-occur with it?
\item \textbf{Stability.} Clinical documentation is highly variable: the same
      finding appears as \texttt{SOB} or \texttt{shortness of breath},
      \texttt{CMO} or \texttt{comfort measures only}, depending on author,
      service and institution. An explanation whose highlights change with
      surface form is unusable at the bedside.
\item \textbf{Plausibility, including direction.} A clinician needs to know not
      only that a phrase mattered but \emph{which way} it pushed. Attention-based
      methods are unsigned by construction and cannot express that a finding such
      as ``ambulating independently, discharged home'' is protective.
\end{enumerate}

To our knowledge no prior work measures all three on the same model and the same
clinical corpus, which means the field tacitly assumes they co-occur. They may
not, and if they do not, the choice of explanation method carries clinical
consequences that are currently invisible.

\medskip\noindent\textbf{Contributions.}
\begin{itemize}
\item A three-axis evaluation framework --- faithfulness, stability,
      plausibility --- applied to six explanation methods on one fine-tuned
      clinical transformer over MIMIC-III, with computational cost as a fourth,
      deployment-oriented axis (Section~\ref{sec:framework}).
\item A random-explanation control on every faithfulness measurement, isolating
      the contribution of the explanation from the confound that masking any
      text degrades a prediction (Section~\ref{sec:faithfulness}).
\item A stability protocol using meaning-preserving clinical paraphrase,
      near-duplicate perturbation, and distractor insertion, quantifying
      robustness to documentation variance (Section~\ref{sec:stability}).
\item A \emph{sign agreement} metric that quantifies what is lost by using an
      unsigned attribution method, turning a qualitative critique of attention
      into a measurable quantity (Section~\ref{sec:plausibility}).
\item A rank-ensemble explainer that attains \TBD\ of SHAP's faithfulness at
      \TBD\ of its cost (Section~\ref{sec:ensemble}).
\item A released, dataset-agnostic implementation, including a synthetic clinical
      corpus with gold rationale spans used to validate the plausibility metric
      itself (Section~\ref{sec:synthetic}).
\end{itemize}

% =====================================================================
\section{Related Work}
% =====================================================================

\subsection{Clinical outcome prediction from notes}
Domain-adapted encoders substantially outperform general-domain BERT on clinical
tasks~\cite{alsentzer2019,lee2020,gu2021}. ClinicalBERT was applied to 30-day
readmission from discharge summaries~\cite{huang2019}, and subsequent work
established a broader benchmark of admission-note outcome prediction
tasks~\cite{vanaken2021}. Our prediction model follows this line and is not
itself a contribution; it exists to provide a realistic, well-trained subject for
the explanation evaluation.

\subsection{Attention as explanation}
Whether attention weights constitute explanation has been contested since
\cite{jain2019} showed attention distributions can be substantially altered
without changing predictions, and \cite{serrano2019} reached similar conclusions
via intermediate-representation erasure. \cite{wiegreffe2019} argued the critique
depends on what one asks of an explanation. Attention
rollout~\cite{abnar2020} addresses a narrower technical objection, propagating
attention across layers while accounting for residual connections, and
\cite{chefer2021} makes the resulting map class-specific by gradient weighting.
We include raw attention, rollout and gradient-weighted rollout as three distinct
conditions rather than treating ``attention'' as one method.

\subsection{Faithfulness evaluation}
ERASER~\cite{deyoung2020} introduced \emph{comprehensiveness} and
\emph{sufficiency} as quantitative faithfulness measures over human-annotated
rationales. Deletion and insertion curves serve a similar role in vision.
Integrated gradients~\cite{sundararajan2017} and SHAP~\cite{lundberg2017} carry
axiomatic guarantees --- completeness and additivity respectively --- that
attention-based methods lack. These metrics are, however, applied predominantly
to general-domain benchmarks; their behaviour on clinical risk prediction, where
notes are long, class imbalance is severe and the deployment stakes differ, is
largely uncharacterised.

\subsection{Stability of explanations}
\cite{ghorbani2019} demonstrated that neural network interpretations are fragile
under imperceptible input perturbation. In clinical text the relevant
perturbation is not adversarial but naturally occurring: the same clinical fact
written by a different author. We are not aware of prior work measuring
explanation invariance to clinical abbreviation and synonym variation, which we
argue is the operationally relevant robustness question in this domain.

\note{TODO before submission: run a targeted search for recent work on joint
faithfulness/stability evaluation in clinical NLP, and on synthetic gold-rationale
benchmarks, to confirm the gap claim in Sec.~I is still accurate.}

% =====================================================================
\section{Data}\label{sec:data}
% =====================================================================

\subsection{MIMIC-III cohort}
We use discharge summaries from MIMIC-III~\cite{johnson2016}. The unit of
analysis is the hospital admission (\texttt{HADM\_ID}). Notes are
\texttt{NOTEEVENTS} rows with \texttt{CATEGORY = `Discharge summary'}; where an
admission has addenda, report and addenda are concatenated in \texttt{CHARTDATE}
order, and rows flagged \texttt{ISERROR} are dropped. The primary outcome is
in-hospital mortality (\texttt{HOSPITAL\_EXPIRE\_FLAG}); 30-day readmission is
reported as a secondary outcome, defined as a non-elective admission for the same
\texttt{SUBJECT\_ID} beginning within 30 days of discharge, with index admissions
ending in death excluded. We exclude patients under 18 at admission, admissions
without a discharge summary, and organ-donor accounts. The resulting cohort
contains \TBD\ admissions with a mortality rate of \TBD.

\subsection{Leakage control}
Two forms of leakage would inflate results and are controlled explicitly.

\textbf{Outcome leakage.} A discharge summary for a patient who died frequently
states the outcome directly (``patient expired'', ``time of death'',
``\texttt{Discharge Disposition: EXPIRED}''). Training on such text is not
prediction. We redact twelve such patterns and report the ablation with and
without redaction (Table~\ref{tab:main}, \note{add ablation row}).

\textbf{Identity leakage.} MIMIC-III patients have repeat admissions. We split on
\texttt{SUBJECT\_ID} with stratification on the outcome, so that no patient
appears in more than one split, and assert both patient disjointness and the
absence of duplicate note text across splits.

\subsection{Preprocessing}
De-identification placeholders are replaced with typed tokens
(\texttt{<DATE>}, \texttt{<PLACE>}, \texttt{<NAME>}) rather than deleted, so that
character offsets used by the explanation payload remain valid. Administrative
header lines are removed; in our corpus this reduced mean note length from
\ph{459} to \ph{341} wordpieces, materially changing what survives truncation.
Abbreviation expansion alters token boundaries and therefore explanations, so it
is disabled by default and reported only as an ablation.

Median note length is \TBD\ wordpieces, with \TBD\% exceeding the 512-token
encoder limit. We use head--tail truncation, retaining the first 35\% and last
65\% of the budget: plain right-truncation discards the discharge disposition,
where goals-of-care language --- among the strongest mortality signals in
narrative notes --- is located.

\subsection{Synthetic corpus for metric validation}\label{sec:synthetic}
Plausibility on real notes requires clinician annotation, which is expensive and
available only at small scale. To validate that our automatic plausibility metric
measures what we claim, we additionally construct a synthetic corpus of
\ph{3{,}000} notes imitating the surface structure of MIMIC-III discharge
summaries, generated from a known process
\begin{equation}
\mathrm{logit} = \beta_0 + \beta_{\text{age}} + \sum_{f \in F} w_f,
\qquad y \sim \mathrm{Bernoulli}(\sigma(\mathrm{logit}))
\end{equation}
over \ph{24} clinical risk factors $F$ with known log-odds weights $w_f$. Because
$y$ is sampled rather than thresholded, Bayes-optimal AUROC is strictly below
1.0. Each injected risk-factor sentence is recorded with its character span,
providing \emph{gold rationales} without annotation. Outcome-neutral distractor
sentences are injected in every section, penalising explainers that attach to
clinical boilerplate. We stress that this corpus validates the metric; all
reported clinical results are on MIMIC-III.

% =====================================================================
\section{Model}\label{sec:model}
% =====================================================================

We fine-tune four encoders for binary classification: Bio\_ClinicalBERT,
BioBERT, PubMedBERT, and BERT-base as a control isolating the benefit of domain
adaptation. Each uses a \texttt{[CLS]} pooled representation followed by dropout
and a linear head.

Class imbalance (\TBD\% positive) is handled with inverse-frequency weighted
cross-entropy; unweighted cross-entropy and focal loss~\cite{lin2017} are
reported as ablations. Model selection uses \textbf{validation AUPRC} rather than
loss or accuracy: at this prevalence, validation loss is dominated by the
majority class and routinely selects a checkpoint that ranks the minority class
worse. The classification head uses a learning rate $10\times$ that of the
encoder. Each configuration is trained across \ph{five} random seeds and reported
as mean $\pm$ standard deviation.

The best model is selected on \emph{validation} AUPRC and evaluated once on the
held-out test set; all explanation experiments use that single checkpoint.

% =====================================================================
\section{Explanation Methods}\label{sec:methods}
% =====================================================================

We compare six methods spanning three families. All return an identical payload
--- per-token scores with character offsets into the source note, aggregated to
whitespace words --- so that the same metrics and the same interface consume all
six without modification.

\subsection{Attention-based}
\textbf{Raw attention} takes the mean over heads of the final layer's
\texttt{[CLS]} attention row. \textbf{Attention rollout}~\cite{abnar2020}
propagates attention across all layers,
\begin{equation}
\hat{A}_l = \mathrm{norm}\!\left(\tfrac{1}{2}\,\overline{A_l} + \tfrac{1}{2}I\right),
\qquad
R = \hat{A}_L \hat{A}_{L-1} \cdots \hat{A}_1,
\end{equation}
where $\overline{A_l}$ is the head-averaged attention at layer $l$ and the
identity term models the residual stream; importance is the \texttt{[CLS]} row of
$R$. Both are unsigned and, in the case of rollout, class-agnostic: the same map
is produced regardless of which class is predicted.
\textbf{Gradient-weighted rollout}~\cite{chefer2021} weights heads by the
gradient of the target logit before fusion, yielding a signed, class-specific map.

\subsection{Perturbation-based}
\textbf{Occlusion} masks each word in turn and records the change in
positive-class log-odds, requiring one forward pass per word.
\textbf{SHAP}~\cite{lundberg2017} is computed with a partition explainer over the
note, exploiting hierarchical text structure for $O(n \log n)$ rather than $2^n$
model evaluations. We compute SHAP values in \emph{log-odds} space: Shapley
additivity holds on a scale where the model is locally linear, and on the
probability scale values cease to sum to the prediction once the sigmoid
saturates --- precisely the high-risk regime of interest.

\subsection{Gradient-based}
\textbf{Integrated gradients}~\cite{sundararajan2017} integrates the gradient
along a path from a baseline $x'$ to the input,
\begin{equation}
\mathrm{IG}_i = (x_i - x'_i)\int_0^1
\frac{\partial F\!\left(x' + \alpha (x - x')\right)}{\partial x_i}\, d\alpha ,
\end{equation}
approximated by Riemann midpoint sum. We use an all-\texttt{[MASK]} baseline with
special tokens preserved rather than the zero embedding, which is not a point in
the model's input distribution for text. Completeness --- attributions summing to
$F(x) - F(x')$ --- provides a built-in correctness check, reported as
convergence error.

% =====================================================================
\section{Evaluation Framework}\label{sec:framework}
% =====================================================================

\subsection{Faithfulness}\label{sec:faithfulness}
Following ERASER~\cite{deyoung2020}, for a prediction $\hat{y}$ with the top-$k$
scoring tokens $r_k$:
\begin{align}
\mathrm{comp} &= p(\hat{y}\mid x) - p(\hat{y}\mid x \setminus r_k), \\
\mathrm{suff} &= p(\hat{y}\mid x) - p(\hat{y}\mid r_k).
\end{align}
Both are averaged over $k \in \{1,5,10,20,50\}\%$ of tokens. Three
implementation choices are load-bearing:

\emph{Removal means masking, not deletion.} Deleting tokens shortens the sequence
and shifts all positions, so part of the measured drop would be a length artefact
rather than evidence removal.

\emph{Every measurement carries a random-explanation control.} Masking any 20\%
of a note degrades confidence to some degree. The quantity that separates methods
is how much \emph{more} the explanation's tokens hurt than a random selection of
the same size, which we report as $\mathrm{comp}_{\Delta\text{rand}}$. Omitting
this control conflates explanation quality with the generic effect of perturbing
text.

\emph{Tokens are ranked by signed score toward the predicted class}, not by
absolute value, which would mix in the strongest evidence \emph{against} the
prediction and understate comprehensiveness.

We additionally report deletion and insertion curves and their areas.

\subsection{Stability}\label{sec:stability}
Each note is re-explained under three meaning-preserving perturbations of
increasing severity: \textbf{near-duplicate} (whitespace and casing only),
\textbf{paraphrase} (clinical abbreviation and synonym substitution), and
\textbf{distractor} (an outcome-irrelevant sentence appended). Agreement is
measured by Spearman correlation over shared vocabulary, top-$k$ Jaccard overlap,
and, for the distractor condition, \emph{distractor mass}: the fraction of total
absolute importance landing on the inserted sentence. We also report the shift in
predicted probability, since a genuine change in the model's output would make a
change in explanation legitimate.

\subsection{Plausibility}\label{sec:plausibility}
On MIMIC-III we conduct a blinded review. For each of \ph{40} test notes, the
top-8 highlighted phrases from each method are presented under neutral labels
that are reshuffled per note, so a rater cannot learn a method's identity.
\ph{Three} clinically-informed raters score clinical sensibility and completeness
on a 1--5 Likert rubric. We report per-method means and Krippendorff's $\alpha$
for ordinal data, chosen over Fleiss' $\kappa$ because it accommodates missing
ratings and respects the ordering of the scale.

On the synthetic corpus, where gold rationales are known, we additionally report
token-level F1 and rationale AUPRC against the gold spans, Spearman correlation
between attribution magnitude and true log-odds weight, and \textbf{sign
agreement}: the fraction of gold rationale tokens whose attributed sign matches
the sign of the true weight. Sign agreement is undefined for unsigned methods,
which is itself the finding --- it quantifies the clinical information those
methods structurally cannot convey.

\subsection{Cost}
We report mean wall-clock seconds and mean model evaluations per note, and plot
faithfulness against cost. For bedside deployment this trade-off, not
faithfulness alone, determines what is usable.

\subsection{Statistical procedure}
Comparisons between models use paired tests, since all models are evaluated on
identical patients: DeLong~\cite{delong1988} for correlated AUROCs and a paired
stratified bootstrap for AUPRC, which has no closed-form variance. Multiple
comparisons against baselines are Holm-corrected. Bootstrap resampling is
stratified by class, as an unstratified bootstrap on a rare outcome produces
folds in which AUROC is undefined.

% =====================================================================
\section{Experimental Setup}
% =====================================================================

Models are fine-tuned for up to \ph{6} epochs with early stopping on validation
AUPRC (patience \ph{2}), batch size \ph{16}, encoder learning rate
\ph{$2\times10^{-5}$}, maximum sequence length \ph{512}, on \ph{a single NVIDIA
A100}. SHAP uses \ph{500} evaluations per note; integrated gradients uses
\ph{64} steps. Faithfulness is computed on \ph{200} test notes and stability on
\ph{150}. Full configuration files, seeds and package versions are released with
the code.

% =====================================================================
\section{Results}
% =====================================================================

\subsection{Predictive performance}
Table~\ref{tab:main} reports held-out test performance. \TBD

\begin{table}[t]
\caption{Held-out test performance on in-hospital mortality, MIMIC-III.
Mean $\pm$ SD over \ph{5} seeds. $p$-values are paired bootstrap against the best
model, Holm-adjusted.}
\label{tab:main}
\centering
\begin{tabular}{lcccc}
\toprule
Model & AUROC & AUPRC & F1 & $p$ \\
\midrule
TF-IDF + LR        & \TBD & \TBD & \TBD & \TBD \\
TF-IDF + XGBoost   & \TBD & \TBD & \TBD & \TBD \\
BERT-base          & \TBD & \TBD & \TBD & \TBD \\
BioBERT            & \TBD & \TBD & \TBD & \TBD \\
PubMedBERT         & \TBD & \TBD & \TBD & \TBD \\
Bio\_ClinicalBERT  & \TBD & \TBD & \TBD & -- \\
\bottomrule
\end{tabular}
\end{table}

\subsection{Faithfulness}
Table~\ref{tab:faith} reports comprehensiveness and sufficiency with the random
control. \TBD

\begin{table}[t]
\caption{Explanation faithfulness on \ph{200} test notes.
$\uparrow$/$\downarrow$ indicate preferred direction.}
\label{tab:faith}
\centering
\begin{tabular}{lcccc}
\toprule
Method & Comp.\,$\uparrow$ & Suff.\,$\downarrow$ & vs.\ rand.\,$\uparrow$ & s/note \\
\midrule
Raw attention       & \TBD & \TBD & \TBD & \TBD \\
Attention rollout   & \TBD & \TBD & \TBD & \TBD \\
Grad-weighted attn. & \TBD & \TBD & \TBD & \TBD \\
Integrated grads    & \TBD & \TBD & \TBD & \TBD \\
Occlusion           & \TBD & \TBD & \TBD & \TBD \\
SHAP                & \TBD & \TBD & \TBD & \TBD \\
\midrule
Rank ensemble (ours)& \TBD & \TBD & \TBD & \TBD \\
\bottomrule
\end{tabular}
\end{table}

\subsection{Stability}
\TBD\ \note{Report Spearman and top-10 Jaccard under paraphrase; highlight
distractor mass.}

\subsection{Plausibility}
\TBD\ \note{Report Likert means, Krippendorff's alpha, and sign agreement.
Sign agreement is undefined for the two unsigned methods -- state this as a
result, not a gap in the table.}

\subsection{Do the axes agree?}
This is the central analysis. We rank the six methods on each axis and compute
rank correlation between axes. \TBD

\begin{figure}[t]
\centering
% \includegraphics[width=\columnwidth]{../results/figures/fig6_cost_vs_faithfulness.png}
\fbox{\parbox[c][3.2cm][c]{0.92\columnwidth}{\centering
\ph{Figure placeholder: faithfulness vs.\ compute cost, log $x$-axis.\\
Generated by \texttt{scripts/07\_figures\_tables.py}}}}
\caption{Faithfulness against computational cost per note. Methods in the
upper-left are preferable for deployment.}
\label{fig:cost}
\end{figure}

\subsection{Rank-ensemble explainer}\label{sec:ensemble}
We combine normalised attention rollout with occlusion by averaging max-abs
normalised scores. \TBD

% =====================================================================
\section{System}
% =====================================================================

The evaluation harness is released together with a working clinician-facing
application: a FastAPI service exposing prediction and explanation endpoints, and
a React/Next.js dashboard rendering the risk score, the note with word-level
highlighting, and a ranked contribution chart, with a toggle between explanation
methods.

Two interface decisions follow directly from our findings. Unsigned methods are
rendered on a single-hue scale with the ``protective'' panel suppressed, since
displaying a low-attention word in a contrasting colour would assert a direction
the method cannot express. The decision threshold is drawn explicitly on the risk
gauge, because for a rare outcome the operating point is far below 0.5 and a
gauge without it invites the reading that below half the dial is safe.

% =====================================================================
\section{Discussion}
% =====================================================================

\TBD

\note{Write after results. If the axes dissociate, the message is: explainer
choice is a clinical decision requiring justification, and papers reporting a
single method without faithfulness evidence are underdetermined. If the axes
converge, the message is: cheap methods suffice, which is a deployment
recommendation. Decide framing from the numbers, never before.}

% =====================================================================
\section{Limitations}
% =====================================================================

Our study covers a single institution's data and a single clinical outcome;
generalisation to other cohorts and outcomes is untested. The plausibility study
uses \ph{three} raters on \ph{40} notes, which is adequate to detect large
differences between methods but not small ones. Sign agreement is measured on the
synthetic corpus, where the true direction of each factor is known by
construction; on MIMIC-III we rely on rater judgement. Attention rollout is
class-agnostic by construction, which limits what an attention-only system can
claim but is a property of the method rather than a defect of our
implementation. Finally, the system is a research prototype and has not been
validated for clinical use.

% =====================================================================
\section{Conclusion}
% =====================================================================

\TBD

% =====================================================================
\section*{Acknowledgment}
% =====================================================================
\ph{Acknowledgements. Note PhysioNet data use agreement here if required.}

% =====================================================================
\begin{thebibliography}{00}

\bibitem{alsentzer2019} E. Alsentzer \emph{et al.}, ``Publicly available clinical BERT embeddings,'' in \emph{Proc. Clinical NLP Workshop}, 2019, pp. 72--78.

\bibitem{lee2020} J. Lee \emph{et al.}, ``BioBERT: a pre-trained biomedical language representation model for biomedical text mining,'' \emph{Bioinformatics}, vol. 36, no. 4, pp. 1234--1240, 2020.

\bibitem{gu2021} Y. Gu \emph{et al.}, ``Domain-specific language model pretraining for biomedical natural language processing,'' \emph{ACM Trans. Comput. Healthcare}, vol. 3, no. 1, pp. 1--23, 2021.

\bibitem{huang2019} K. Huang, J. Altosaar, and R. Ranganath, ``ClinicalBERT: Modeling clinical notes and predicting hospital readmission,'' \emph{arXiv:1904.05342}, 2019.

\bibitem{vanaken2021} B. van Aken \emph{et al.}, ``Clinical outcome prediction from admission notes using self-supervised knowledge integration,'' in \emph{Proc. EACL}, 2021, pp. 881--893.

\bibitem{jain2019} S. Jain and B. C. Wallace, ``Attention is not explanation,'' in \emph{Proc. NAACL-HLT}, 2019, pp. 3543--3556.

\bibitem{wiegreffe2019} S. Wiegreffe and Y. Pinter, ``Attention is not not explanation,'' in \emph{Proc. EMNLP-IJCNLP}, 2019, pp. 11--20.

\bibitem{serrano2019} S. Serrano and N. A. Smith, ``Is attention interpretable?,'' in \emph{Proc. ACL}, 2019, pp. 2931--2951.

\bibitem{abnar2020} S. Abnar and W. Zuidema, ``Quantifying attention flow in transformers,'' in \emph{Proc. ACL}, 2020, pp. 4190--4197.

\bibitem{chefer2021} H. Chefer, S. Gur, and L. Wolf, ``Transformer interpretability beyond attention visualization,'' in \emph{Proc. CVPR}, 2021, pp. 782--791.

\bibitem{deyoung2020} J. DeYoung \emph{et al.}, ``ERASER: A benchmark to evaluate rationalized NLP models,'' in \emph{Proc. ACL}, 2020, pp. 4443--4458.

\bibitem{lundberg2017} S. M. Lundberg and S.-I. Lee, ``A unified approach to interpreting model predictions,'' in \emph{Proc. NeurIPS}, 2017, pp. 4765--4774.

\bibitem{sundararajan2017} M. Sundararajan, A. Taly, and Q. Yan, ``Axiomatic attribution for deep networks,'' in \emph{Proc. ICML}, 2017, pp. 3319--3328.

\bibitem{ghorbani2019} A. Ghorbani, A. Abid, and J. Zou, ``Interpretation of neural networks is fragile,'' in \emph{Proc. AAAI}, 2019, pp. 3681--3688.

\bibitem{johnson2016} A. E. W. Johnson \emph{et al.}, ``MIMIC-III, a freely accessible critical care database,'' \emph{Scientific Data}, vol. 3, 160035, 2016.

\bibitem{delong1988} E. R. DeLong, D. M. DeLong, and D. L. Clarke-Pearson, ``Comparing the areas under two or more correlated receiver operating characteristic curves: a nonparametric approach,'' \emph{Biometrics}, vol. 44, no. 3, pp. 837--845, 1988.

\bibitem{lin2017} T.-Y. Lin \emph{et al.}, ``Focal loss for dense object detection,'' in \emph{Proc. ICCV}, 2017, pp. 2980--2988.

\end{thebibliography}

\end{document}
